\ifdefined\gkver\else\def\gkver{student}\fi
\documentclass[\gkver]{gaokaozhenti}
\begin{document}

\chapter{2022年海南}

\begin{enumerate}
\item
%% number: 1
%% paperName: 2022年海南高考第 1 题 3 分
%% typeId: 1
%% type: 单选题
%% chapter: 运动和力
%% point: 动量和冲量
%% method: 无
%% score: 3
%% degree: 950
%% duplicateId: 0
%% body:
北京 2022 年冬奥会，中国短道速滑队获得混合接力冠军。如图，交接棒的过程中甲运动员用力推乙运动员，设此过程中甲对乙的作用力为 $F_{1}$，乙对甲的作用力为 $F_{2}$，则\xzanswer{A}

\onepicture[w=5.60cm]{figs/01.jpg}

\fourchoices[answer=A]
{$F_{1}$ 和 $F_{2}$ 总是大小相等，方向相反}
{$F_{1}$ 和 $F_{2}$ 总是大小相等，方向相同}
{$F_{1}$ 的冲量大小比 $F_{2}$ 的大}
{$F_{1}$ 的冲量大小比 $F_{2}$ 的小}

%% answer: A
%% memo:
\memoanswer{
	根据题意可知 $F_{1}$ 和 $F_{2}$ 是相互作用力，根据牛顿第三定律可知 $F_{1}$ 和 $F_{2}$ 等大反向、具有同时性；根据冲量定义式 $I = Ft$ 可知 $F_{1}$ 和 $F_{2}$ 的冲量大小相等，方向相反。

	故选 A。
}

\item
%% number: 2
%% paperName: 2022年海南高考第 2 题 3 分
%% typeId: 1
%% type: 单选题
%% chapter: 原子和原子核
%% point: 天然放射性
%% method: 无
%% score: 3
%% degree: 950
%% duplicateId: 0
%% body:
下列四个核反应方程中，属于 $\beta$ 衰变的是\xzanswer{C}

\fourchoices[answer=C]
{\ce{^{238}_{92} U -> ^{234}_{90} Th + ^{4}_{2} He}}
{\ce{^{14}_{7} N + ^{4}_{2} He -> ^{17}_{8} O + ^{1}_{1} H}}
{\ce{^{234}_{90} Th -> ^{234}_{91} Pa + ^{0}_{-1} e}}
{\ce{^{235}_{92} U + ^{1}_{0} n -> ^{144}_{56} Ba + ^{89}_{36} Kr + 3^{1}_{0} n}}

%% answer: C
%% memo:
\memoanswer{
	A．该反应属于 $\alpha$ 衰变，放出了氦核（\ce{^{4}_{2}He}），A 错误；

	B．该反应是卢瑟福发现质子（\ce{^{1}_{1}H}）的核反应方程，B 错误；

	C．该反应属于 $\beta$ 衰变，放出了电子（\ce{^{0}_{-1}e}），C 正确；

	D．该反应是重核裂变的核反应方程，D 错误。

	故选 C。
}

\item
%% number: 3
%% paperName: 2022年海南高考第 3 题 3 分
%% typeId: 1
%% type: 单选题
%% chapter: 光学
%% point: 光的折射
%% method: 无
%% score: 3
%% degree: 750
%% duplicateId: 0
%% body:
空气中有一透明材料制成的中空管，其横截面内外圆半径之比 $R_{1}:R_{2} = 1:3$。一束单色光以 $i = 30^\circ$ 的入射角平行于横截面从 A 点入射，经折射后与内圆相切，如图所示。该材料的折射率为\xzanswer{B}

\onepicture[w=4.60cm]{figs/03.png}

\fourchoices[answer=B]
{$\frac{4}{3}$}
{$\frac{3}{2}$}
{$\frac{\sqrt{10}}{2}$}
{$\sqrt{3}$}

%% answer: B
%% memo:
\memoanswer{
	如图所示，设光经过 A 点后的折射角为 $r$，$r$ 的正弦值 $\sin r = \frac{R_{1}}{R_{2}} = \frac{1}{3}$，由折射定律可得该透明材料的折射率 $n = \frac{\sin 30^\circ}{\sin r} = 1.5$，故 B 正确。

	\onepicture[w=4.20cm]{figs/03_an.png}
}

\item
%% number: 4
%% paperName: 2022年海南高考第 4 题 3 分
%% typeId: 1
%% type: 单选题
%% chapter: 机械振动
%% point: 振动图象
%% method: 无
%% score: 3
%% degree: 650
%% duplicateId: 0
%% body:
甲、乙两摆在同一地点的振动图像如图所示，两单摆的摆长分别为 $l_{\text{甲}}$、$l_{\text{乙}}$，则 $\frac{l_{\text{甲}}}{l_{\text{乙}}}$ 为\xzanswer{C}

\onepicture[w=5.60cm]{figs/04.png}

\fourchoices[answer=C]
{$\frac{2}{3}$}
{$\frac{3}{2}$}
{$\frac{4}{9}$}
{$\frac{9}{4}$}

%% answer: C
%% memo:
\memoanswer{
	由振动图像可知甲、乙两个单摆周期之比为 $T_{\text{甲}}:T_{\text{乙}} = 0.8:1.2 = 2:3$。根据单摆周期公式可得 $l \propto T^{2}$，则甲、乙两个单摆的摆长之比为 $l_{\text{甲}}:l_{\text{乙}} = T_{\text{甲}}^{2}:T_{\text{乙}}^{2} = 4:9$。

	故选 C。
}

\item
%% number: 5
%% paperName: 2022年海南高考第 5 题 3 分
%% typeId: 1
%% type: 单选题
%% chapter: 交流电
%% point: 交流电
%% method: 无
%% score: 3
%% degree: 750
%% duplicateId: 0
%% body:
如图，矩形导线框在匀强磁场中绕垂直于磁场的轴 $OO'$ 以角速度 $\omega$ 匀速转动。线框面积为 $S$，匝数为 $N$，磁感应强度大小为 $B$，$t = 0$ 时线框平面与磁场方向平行。下列四幅图中能正确表示线框上产生的感应电动势 $e$ 随时间 $t$ 变化关系的是\xzanswer{A}

\onepicture[w=3.60cm]{figs/05.png}

\fourchoices[answer=A, ispicture=true, w=3.40cm]
{figs/05a.png}
{figs/05b.png}
{figs/05c.png}
{figs/05d.png}

%% answer: A
%% memo:
\memoanswer{
	图示位置线框处于与中性面垂直的平面，竖直长边垂直切割磁感线，此时产生的感应电动势最大为 $E_{m} = NBS\omega$。根据正弦式交变电流的表达式可知感应电动势随时间的变化关系为 $e = NBS\omega\cos\omega t$。

	故选 A。
}

\item
%% number: 6
%% paperName: 2022年海南高考第 6 题 3 分
%% typeId: 1
%% type: 单选题
%% chapter: 磁场
%% point: 磁场的叠加
%% method: 无
%% score: 3
%% degree: 550
%% duplicateId: 0
%% body:
如图，四根与纸面垂直的固定长直细导线，横截面分别位于正方形的四个顶点，通有大小相等的电流，其中导线 $L_{1}$、$L_{2}$、$L_{3}$ 中的电流方向垂直纸面向里，导线 $L_{4}$ 中的电流方向垂直纸面向外。$O$ 点为正方形的中心，已知每根导线中的电流在 $O$ 点产生的磁感应强度大小为 $B$，则 $O$ 点的磁感应强度\xzanswer{A}

\onepicture[w=4.20cm]{figs/06.png}

\fourchoices[answer=A]
{方向由 $O$ 点沿纸面指向 $L_{1}$，大小为 $2B$}
{方向由 $O$ 点沿纸面指向 $L_{2}$，大小为 $B$}
{方向由 $O$ 点沿纸面指向 $L_{3}$，大小为 $2B$}
{方向由 $O$ 点沿纸面指向 $L_{4}$，大小为 $B$}

%% answer: A
%% memo:
\memoanswer{
	$L_{1}$ 和 $L_{3}$ 导线在 $O$ 点产生的磁感应强度等大反向，$L_{2}$ 和 $L_{4}$ 导线在 $O$ 点产生的磁感应强度等大同向，方向指向 $L_{1}$，所以 $O$ 点的磁感应强度大小为 $2B$，方向指向 $L_{1}$，故 A 正确。
}

\item
%% number: 7
%% paperName: 2022年海南高考第 7 题 3 分
%% typeId: 1
%% type: 单选题
%% chapter: 磁场
%% point: 带电粒子在磁场中的运动
%% method: 无
%% score: 3
%% degree: 650
%% duplicateId: 0
%% body:
如图，两块圆弧形金属板间存在方向指向圆心 $O$ 的电场，与 $O$ 点等距处电场强度大小相等。一束正离子流沿纸面垂直电场方向以相同速度射入电场，其中一部分离子能沿着某一等势面通过电场，然后进入方向垂直纸面向外的匀强磁场，则这部分离子中，沿相同轨迹通过磁场的离子都具有相同的\xzanswer{C}

\onepicture[w=5.00cm]{figs/07.png}

\fourchoices[answer=C]
{质量}
{电荷量}
{比荷}
{动能}

%% answer: C
%% memo:
\memoanswer{
	粒子在辐射电场中以速度 $v$ 做匀速圆周运动，电场力完全提供向心力，根据牛顿第二定律可知

	$qE = m\frac{v^{2}}{r}$

	解得 $r = \frac{m}{q} \cdot \frac{v^{2}}{E}$

	粒子在匀强磁场中 $qvB = m\frac{v^{2}}{r}$

	解得 $r' = \frac{m}{q} \cdot \frac{v}{B}$

	粒子不同场中的轨迹相同，即粒子在不同场中转动半径相同，所以这些粒子具有相同的速度 $v$ 和比荷 $\frac{q}{m}$。

	故选 C。
}

\item
%% number: 8
%% paperName: 2022年海南高考第 8 题 3 分
%% typeId: 1
%% type: 单选题
%% chapter: 力物体的平衡
%% point: 三力平衡
%% method: 无
%% score: 3
%% degree: 550
%% duplicateId: 0
%% body:
具有拱券结构的赵州桥展示了我国古代人民高超的造桥技术。如图，某同学用 6 块形状相同的楔形块搭成一个半圆形的拱券结构桥梁模型，其中 1 号与 6 号楔形块固定在水平面上，2 号与 5 号楔形块的质量均为 $m$，3 号与 4 号楔形块的质量均为 $m'$。若不计楔形块之间的摩擦力，则 $\frac{m}{m'}$ 为\xzanswer{D}

\onepicture[w=5.40cm]{figs/08.png}

\fourchoices[answer=D]
{$\frac{\sqrt{3}}{2}$}
{1}
{$\sqrt{3}$}
{2}

%% answer: D
%% memo:
\memoanswer{
	六块形状完全相同的石块围成半圆对应的圆心角为 $180^\circ$，每块石块对应的圆心角为 $30^\circ$，对第 3 块石块受力分析如图~\ref{2022海南08a}所示。结合力的合成可知

	$\tan 60^\circ = \frac{F_{4}}{m'g}$

	对第 2 块和第 3 块石块整体受力分析如图~\ref{2022海南08b}所示

	$\tan 30^\circ = \frac{F_{4}}{(m + m')g}$

	解得

	$\frac{m}{m'} = 2$

	故选 D。

	\twopicture[labelA=2022海南08a, labelB=2022海南08b, numA=甲, numB=乙, wA=3.60cm, wB=3.60cm]{figs/08_an1.png}{figs/08_an2.png}
}

\item
%% number: 9
%% paperName: 2022年海南高考第 9 题 4 分
%% typeId: 2
%% type: 多选题
%% chapter: 电场
%% point: 电场线
%% method: 无
%% score: 4
%% degree: 750
%% duplicateId: 0
%% body:
某带电体周围的电场线和等势面如图所示，图中 A、B 两点的电场强度大小分别为 $E_{A}$、$E_{B}$，电势分别为 $\varphi_{A}$、$\varphi_{B}$，则\xzanswer{BD}

\onepicture[w=5.00cm]{figs/09.png}

\fourchoices[answer=BD]
{$E_{A} > E_{B}$}
{$E_{A} < E_{B}$}
{$\varphi_{A} > \varphi_{B}$}
{$\varphi_{A} < \varphi_{B}$}

%% answer: BD
%% memo:
\memoanswer{
	A、B．根据电场线的疏密程度表示电场的强弱，电场线越密，电场强度越大，则 B 点的电场强度较大，即 $E_{A} < E_{B}$。故 A 错误，B 正确；

	C、D．根据电场线与等势面垂直，且电势较高的等势面指向电势较低的等势面，所以 B 点的电势较大，即 $\varphi_{A} < \varphi_{B}$。故 C 错误，D 正确。

	故选 BD。
}

\item
%% number: 10
%% paperName: 2022年海南高考第 10 题 4 分
%% typeId: 2
%% type: 多选题
%% chapter: 曲线运动
%% point: 万有引力定律
%% method: 无
%% score: 4
%% degree: 550
%% duplicateId: 0
%% body:
火星与地球的质量之比为 $a$，半径之比为 $b$。设火星、地球的第一宇宙速度分别为 $v_{\text{火}}$、$v_{\text{地}}$，表面的自由落体加速度分别为 $g_{\text{火}}$、$g_{\text{地}}$，则\xzanswer{AC}

\fourchoices[answer=AC]
{$\frac{v_{\text{火}}}{v_{\text{地}}} = \sqrt{\frac{a}{b}}$}
{$\frac{v_{\text{火}}}{v_{\text{地}}} = \sqrt{\frac{b}{a}}$}
{$\frac{g_{\text{火}}}{g_{\text{地}}} = \frac{a}{b^{2}}$}
{$\frac{g_{\text{火}}}{g_{\text{地}}} = \frac{a}{b}$}

%% answer: AC
%% memo:
\memoanswer{
	由 $G\frac{Mm}{R^{2}} = m\frac{v^{2}}{R}$ 结合 $gR^{2} = GM$ 可得 $v = \sqrt{gR}$，知 $\frac{v_{\text{火}}}{v_{\text{地}}} = \sqrt{\frac{a}{b^{2}} \cdot b} = \sqrt{\frac{a}{b}}$。

	由 $G\frac{Mm}{R^{2}} = mg$ 可得 $g = \frac{GM}{R^{2}}$，知 $\frac{g_{\text{火}}}{g_{\text{地}}} = \frac{a}{b^{2}}$。

	故 AC 正确，BD 错误。

	故选 AC。
}

\item
%% number: 11
%% paperName: 2022年海南高考第 11 题 4 分
%% typeId: 2
%% type: 多选题
%% chapter: 原子和原子核
%% point: 玻尔模型
%% method: 无
%% score: 4
%% degree: 650
%% duplicateId: 0
%% body:
大量氢原子处于 $n = 4$ 能级，会自发地辐射出多种频率的光，下列说法正确的是\xzanswer{BC}

\fourchoices[answer=BC]
{向基态跃迁时氢原子吸收能量}
{能辐射出 6 种频率的光}
{从 $n = 4$ 能级跃迁到 $n = 1$ 能级时辐射的光频率最大}
{从 $n = 4$ 能级跃迁到 $n = 3$ 能级时辐射的光频率最大}

%% answer: BC
%% memo:
\memoanswer{
	A．高能级向低能级跃迁向外放出能量，以光子形式释放出去，故 A 错误；

	B．最多能放不同频率光子的种数为 $C_{4}^{2} = 6$。故 B 正确；

	C、D．从最高能级向最低能级跃迁释放的光子能量最大，对应的频率最大，波长最小，则 $n = 4$ 向 $n = 1$ 跃迁发出的光子频率最大，故 C 正确，D 错误。

	故选 BC。
}

\item
%% number: 12
%% paperName: 2022年海南高考第 12 题 4 分
%% typeId: 2
%% type: 多选题
%% chapter: 直线运动
%% point: 运动图象
%% method: 无
%% score: 4
%% degree: 450
%% duplicateId: 0
%% body:
某同学从 5 楼乘电梯下降到 1 楼，在此过程中他用手机内置的加速度传感器测得电梯运行的加速度 $a$ 随时间 $t$ 的变化关系如图所示，则\xzanswer{AD}

\onepicture[w=5.60cm]{figs/12.png}

\fourchoices[answer=AD]
{$t = 1 \Us$ 时该同学处于失重状态}
{$t = 10 \Us$ 时该同学处于失重状态}
{此过程中电梯运行中最大的速率约为 $0.56 \Ums$}
{此过程中电梯运行中最大的速率约为 $1.8 \Ums$}

%% answer: AD
%% memo:
\memoanswer{
	因为电梯向下运动，则最初电梯加速度方向竖直向下，由题意知 $a < 0$ 表示加速度方向竖直向下。$t = 1 \Us$ 时，$a < 0$，即加速度方向竖直向下，电梯对该同学的支持力小于重力，该同学处于失重状态，故 A 正确；

	$t = 10 \Us$ 时，$a > 0$，加速度方向竖直向上，电梯对该同学的支持力大于重力，该同学处于超重状态，故 B 错误；

	$a-t$ 图像中图线与 $t$ 轴围成的面积表示速度的变化量，开始时速率为零，则运动中最大速率约为 $v_{m} = \frac{1.5 + 4.5}{2} \times 0.6 \Ums = 1.8 \Ums$。故 C 错误 D 正确。
}

\item
%% number: 13
%% paperName: 2022年海南高考第 13 题 4 分
%% typeId: 2
%% type: 多选题
%% chapter: 电场
%% point: 带电粒子的直线运动
%% method: 无
%% score: 4
%% degree: 350
%% duplicateId: 0
%% body:
如图~\ref{2022海南13a}，物块 A 放在绝缘水平面上的 $O$ 点，用跨过轻质定滑轮的绝缘轻绳与物块 B 连接，B 离滑轮足够远。空间存在水平向左的匀强电场，电场强度大小为 $4 \times 10^{5} \UVm$。A、B 的质量均为 $1 \Ukg$，A 的电荷量为 $3 \times 10^{-5} \UC$，B 不带电。运动过程中 A 与 $O$ 点的距离设为 $x$，A 与水平面间的动摩擦因数 $\mu$ 与 $x$ 的关系如图~\ref{2022海南13b}。规定 $O$ 点的电势为零，重力加速度大小取 $g = 10 \Umsq$。将 A、B 由静止释放，则物块 A 的\xzanswer{ACD}

\twopicture[labelA=2022海南13a, labelB=2022海南13b, numA=甲, numB=乙, wA=6.20cm, wB=4.60cm]{figs/13a.png}{figs/13b.png}

\fourchoices[answer=ACD]
{最大速度为 $1 \Ums$}
{最大位移为 $1 \Um$}
{当速度为 $0.6 \Ums$ 时，电势能可能是 $-2.4 \UJ$}
{当速度为 $0.6 \Ums$ 时，绳中拉力大小可能是 $9.2 \UN$}

%% answer: ACD
%% memo:
\memoanswer{
	A、B．由题知 $f = \mu mg = 2x$

	设 A 向左移动 $x$ 后速度为零，对 A、B 系统有

	$qEx - mgx - \frac{1}{2}fx = 0$

	（此处 $fx$ 前面的 $\frac{1}{2}$ 是因为摩擦力是变力，其做功可以用平均力），可得 $x = 2 \Um$

	A 向左运动是先加速后减速，当 $x = 2 \Um$ 时，摩擦力变成静摩擦力，并反向，系统受力平衡，最后静止。设 A 向左运动 $x'$ 后速度为 $v$，对系统则有

	$qEx' - mgx' - \frac{1}{2}fx' = \frac{1}{2} \cdot 2mv^{2}$

	得 $mv^{2} = -(x' - 1)^{2} + 1$

	即：当 $x' = 1 \Um$ 时，$v$ 最大为 $1 \Ums$，故 A 正确，B 错误；

	C．当 $v = 0.6 \Ums$ 时，可得 $x = 0.2 \Um$ 或 $1.8 \Um$

	当 $x = 0.2 \Um$ 时，电场力做功 $qEx = 2.4 \UJ$

	则电势能减小 $2.4 \UJ$，由于 $E_{pO} = 0$，则电势能为 $-2.4 \UJ$，当 $x = 1.8 \Um$ 时

	$E_{pO} = -21.6 \UJ$

	故 C 正确；

	D．根据牛顿第二定律 $qE - f - mg = ma$

	当 $x = 0.2 \Um$ 时，系统加速度 $a = 0.8 \Umsq$

	对 B 有 $T - mg = ma$，得 $T = 10.8 \UN$。

	当 $x = 1.8 \Um$ 时，系统加速度 $a = -0.8 \Umsq$

	对 B 分析可得 $T = 9.2 \UN$。

	故 D 正确。

	故选 ACD。
}

\item
%% number: 14
%% paperName: 2022年海南高考第 14 题 4 分
%% typeId: 4
%% type: 实验题
%% chapter: 光学
%% point: 双缝干涉测量波长
%% method: 无
%% score: 4
%% degree: 850
%% duplicateId: 0
%% body:
用双缝干涉测量某单色光的波长。

\begin{enumerate}
	\item
	如图所示的实验装置中，a、b 分别为\tkanswer{A}（填“A”或“B”）。

	\twochoices[answer=A]
	{单缝、双缝}
	{双缝、单缝}
	\item
	测得第 1 条亮条纹中心到第 6 条亮条纹中心的距离为 $x$，双缝到毛玻璃屏的距离为 $L$，已知双缝间的距离为 $d$，则该单色光的波长 $\lambda =$ \tkanswer{$\frac{dx}{5L}$}（用 $x$、$d$、$L$ 表示）。
\end{enumerate}

\onepicture[w=8.00cm, align=right]{figs/14.png}

%% answer:
\jdanswer{
\begin{enumerate}
	\item
	A

	\item
	$\frac{dx}{5L}$
\end{enumerate}
}
%% memo:
\memoanswer{
	（1）由双缝干涉原理可知，先用滤光片得到单色光，然后用单缝得到细长的光源，最后用双缝得到两束相干光，故 a、b 分别为单缝和双缝。

	故选 A。

	（2）第 1 条亮纹到第 6 条亮纹间距是 $x$，则相邻亮条纹间距为

	$\Delta x = \frac{x}{5}$

	根据 $\Delta x = \frac{L\lambda}{d}$ 可得光的波长是

	$\lambda = \frac{\Delta x d}{L} = \frac{dx}{5L}$
}

\item
%% number: 15
%% paperName: 2022年海南高考第 15 题 8 分
%% typeId: 4
%% type: 实验题
%% chapter: 曲线运动
%% point: 研究平抛运动实验
%% method: 无
%% score: 8
%% degree: 750
%% duplicateId: 0
%% body:
某学习小组探究平抛运动的特点。

\begin{enumerate}
	\item
	采用如图~\ref{2022海南15a}所示装置探究平抛运动竖直分运动的特点。用小锤击打弹性金属片后，A 球沿水平方向抛出，做平抛运动；同时 B 球被释放，自由下落，做自由落体运动。实验发现两球同时落地。分别改变小球距地面的高度和小锤击打的力度，多次重复实验，发现两球仍同时落地。根据该实验现象，可以得到的实验结论是\tkanswer{作平抛运动的物体，在竖直方向上是自由落体运动}。
	\item
	探究平抛运动水平分运动的特点时，得到小球平抛运动的轨迹如图~\ref{2022海南15b}所示，其中 $O$ 为抛出点，$a$、$b$、$c$ 是轨迹上选取的三个点，$O$ 与 $a$、$a$ 与 $b$、$b$ 与 $c$ 之间的竖直距离分别为 $h$、$3h$、$5h$，则小球从 $O$ 到 $a$、$a$ 到 $b$、$b$ 到 $c$ 的运动时间\tkanswer{相等}（填“相等”或“不相等”）；又测得 $O$ 与 $a$、$a$ 与 $b$、$b$ 与 $c$ 之间的水平距离相等，均为 $x$，则可分析得出平抛运动在水平方向的分运动是匀速直线运动，小球平抛运动的初速度为\tkanswer{$x\sqrt{\frac{g}{2h}}$}（用 $h$、$x$ 和重力加速度 $g$ 表示）。
\end{enumerate}

\twopicture[labelA=2022海南15a, labelB=2022海南15b, numA=甲, numB=乙, wA=3.00cm, wB=3.60cm, align=right]{\includesvg{figs/15a}}{figs/15b.png}

%% answer:
\jdanswer{
\begin{enumerate}
	\item
	作平抛运动的物体，在竖直方向上是自由落体运动

	\item
	相等；$x\sqrt{\frac{g}{2h}}$
\end{enumerate}
}
%% memo:
\memoanswer{
	（1）经过多次实验发现两个小球总是同时落地，则得到的结论是：作平抛运动的物体，在竖直方向上是自由落体运动；

	（2）在水平方向是匀速运动，由图可知，从 $O \to a$，$a \to b$，$b \to c$ 水平方向位移相等，运动相同的距离所用时间相等；

	在竖直方向上，连续相等时间 $t$ 内的下落高度之比为 $1:3:5$，可知 $O$ 点的竖直速度为 0，得 $t = \sqrt{\frac{2h}{g}}$，而水平方向上有 $x = vt$，所以 $v = x\sqrt{\frac{g}{2h}}$。
}

\item
%% number: 16
%% paperName: 2022年海南高考第 16 题 8 分
%% typeId: 4
%% type: 实验题
%% chapter: 电路
%% point: 电表改装
%% method: 无
%% score: 8
%% degree: 550
%% duplicateId: 0
%% body:
某同学将一个满偏电流 $I_{g} = 50 \UmuA$ 的表头 G 改装成量程为 $1 \UmA$ 的电流表并校准。

\begin{enumerate}
	\item
	用图~\ref{2022海南16a}所示的电路测量表头 G 的内阻 $r_{g}$，连接线路，闭合开关 S，调节滑动变阻器 $R$，当电流表 A 的示数为 $84 \UmA$ 时，电流表 G 的示数如图~\ref{2022海南16b}所示，该读数为\tkanswer{30.0} $\UmuA$。
	\item
	已知电流表 A 的内阻 $r_{A} = 1.0 \UO$，则表头 G 的内阻 $r_{g} =$ \tkanswer{2800} $\UO$。经计算，用一个阻值为 $R_{1}$ 的电阻与表头 G 并联，将表头 G 改装成电流表。
	\item
	将改装后的电流表与标准电流表串联进行校准。当通过表头 G 的电流为 $\frac{4}{5}I_{g}$ 时，标准电流表 A 的示数为 $0.76 \UmA$，则改装后的电流表量程实际为\tkanswer{0.95} $\UmA$。
	\item
	如果因为测量不准确导致改装后的电流表量程不是 $1 \UmA$，为使改装后的电流表量程变为 $1 \UmA$，只需在 $R_{1}$ 两端并联一个阻值 $R_{2} =$ \tkanswer{18}$R_{1}$ 的电阻即可。
\end{enumerate}

\twopicture[labelA=2022海南16a, labelB=2022海南16b, numA=甲, numB=乙, wA=5.00cm, wB=5.60cm, align=right]{figs/16a.png}{figs/16b.png}

%% answer:
\jdanswer{
\begin{enumerate}
	\item
	$30.0 \UmuA$

	\item
	$2800 \UO$

	\item
	$0.95 \UmA$

	\item
	18
\end{enumerate}
}
%% memo:
\memoanswer{
	（1）电流表 G 的满偏电流 $I_{g} = 50 \UmuA$，则如图乙所示电流表 G 的示数为 $30.0 \UmuA$。

	（2）根据并联电路电压相等可得 $84 \times 10^{-3} \UA \times 1.0 \UO = 30 \times 10^{-6} \UA \times r_{g}$，可得电流表 G 的内阻 $r_{g} = 2800 \UO$。

	（3）分析改装后的电表，由改装电表原理知，当流过表头的电流为 $\frac{4}{5}I_{g}$ 时，则流过电阻 $R_{1}$ 的电流为 $I_{R_{1}} = 0.76 \UmA - \frac{4}{5}I_{g}$，由并联电路电压相等得 $I_{R_{1}} \cdot R_{1} = \frac{4}{5}I_{g} \cdot r_{g}$，解得 $R_{1} = \frac{1}{18}r_{g}$，则改装电流表的实际量程为 $I_{1} = I_{g} + \frac{I_{g}r_{g}}{R_{1}} = 0.95 \UmA$。

	（4）若把并联 $R_{1}$ 的电流表 G 再改装成量程为 $1 \UmA$ 的电流表，需要再并联电阻 $R_{2}$，$R_{1}$ 与 $R_{2}$ 并联的总电阻为 $R_{12}$，由并联电路特点得 $\frac{r_{g}}{R_{12}} = \frac{1000 \UmuA - 50 \UmuA}{50 \UmuA} = \frac{19}{1}$，即 $R_{12} = \frac{1}{19}r_{g}$，又 $R_{12} = \frac{R_{1}R_{2}}{R_{1} + R_{2}}$，解得 $R_{2} = r_{g}$，由（3）知 $R_{1} = \frac{1}{18}r_{g}$，故 $R_{2} = 18R_{1}$。
}

\item
%% number: 17
%% paperName: 2022年海南高考第 17 题 10 分
%% typeId: 6
%% type: 计算题
%% chapter: 气体性质
%% point: 气体多过程
%% method: 无
%% score: 10
%% degree: 650
%% duplicateId: 0
%% body:
如图~\ref{2022海南17a}，粗细均匀、足够长的细玻璃管水平放置，管内用长 $h = 19 \Ucm$ 的水银柱封闭一定质量的理想气体，气柱长 $l_{1} = 25 \Ucm$，温度 $T_{1} = 300 \UK$，大气压强 $p_{0} = 76 \UcmHg$。保持气体温度不变，将玻璃管顺时针缓慢转至开口竖直向上，如图~\ref{2022海南17b}。

\begin{enumerate}
	\item
	判断此过程气体是吸热还是放热；
	\item
	求玻璃管开口竖直向上时气柱的长度 $l_{2}$；
	\item
	保持玻璃管开口竖直向上，若封闭气体缓慢升温至 $T_{2} = 360 \UK$，求此时气柱的长度 $l_{3}$。
\end{enumerate}

\twopicture[labelA=2022海南17a, labelB=2022海南17b, numA=甲, numB=乙, wA=5.00cm, wB=1.80cm, align=right]{figs/17a.png}{figs/17b.png}

%% answer:
\jdanswer{
\begin{enumerate}
	\item
	放热

	\item
	$20 \Ucm$

	\item
	$24 \Ucm$
\end{enumerate}
}
%% memo:
\memoanswer{
	（1）以封闭气体为研究对象，气体做等温变化，内能不变，$\Delta U = 0$，但外界对气体做正功，$W > 0$，由热力学第一定律表达式 $\Delta U = W + Q$ 可知，$Q < 0$，气体对外放热。

	（2）设玻璃管横截面积为 $S$，玻璃管水平时封闭气体压强

	$p_{1} = 76 \UcmHg$

	体积 $V_{1} = l_{1}S$

	玻璃管竖起来后压强为

	$p_{2} = 19 \UcmHg + 76 \UcmHg = 95 \UcmHg$

	体积 $V_{2} = l_{2}S$

	封闭气体发生等温变化，由玻意耳定律得

	$p_{1}V_{1} = p_{2}V_{2}$，解得 $l_{2} = 20 \Ucm$。

	（3）保持玻璃管竖直向上，升温过程中气体发生等压变化，由盖-吕萨克定律得

	$\frac{V_{2}}{T_{1}} = \frac{V_{3}}{T_{2}}$

	其中 $T_{1} = 300 \UK$，$T_{2} = 360 \UK$，$V_{2} = l_{2}S$，$V_{3} = l_{3}S$

	解得 $l_{3} = 24 \Ucm$。
}

\item
%% number: 18
%% paperName: 2022年海南高考第 18 题 12 分
%% typeId: 6
%% type: 计算题
%% chapter: 运动和力
%% point: 动量和能量
%% method: 无
%% score: 12
%% degree: 550
%% duplicateId: 0
%% body:
如图，竖直面内一倾角可调的轨道与光滑水平面平滑相连。当轨道倾角为 $30^\circ$ 时，物块 A 恰好能沿轨道匀速下滑。小球 B 用轻绳悬挂于 O 点，恰与水平面接触而无挤压。将轨道倾角调为 $60^\circ$，物块 A 从轨道上距水平面高为 $h$ 处由静止开始运动，一段时间后与小球 B 发生对心弹性碰撞，碰撞后小球 B 恰好能做完整的圆周运动。已知物块 A 的质量是小球 B 的 3 倍，重力加速度为 $g$，忽略空气阻力，求：

\begin{enumerate}
	\item
	物块 A 与轨道间的动摩擦因数；
	\item
	碰撞后瞬间小球 B 的速度大小；
	\item
	轻绳的长度。
\end{enumerate}

\onepicture[w=7.00cm, align=right]{figs/18.png}

%% answer:
\jdanswer{
\begin{enumerate}
	\item
	$\mu = \frac{\sqrt{3}}{3}$

	\item
	$\sqrt{3gh}$

	\item
	$0.6h$
\end{enumerate}
}
%% memo:
\memoanswer{
	（1）倾角为 $30^\circ$ 时匀速运动，根据平衡条件有

	$mg\sin 30^\circ = \mu mg\cos 30^\circ$

	得 $\mu = \frac{\sqrt{3}}{3}$

	（2）（3）A 从高为 $h$ 的地方滑下后速度为 $v_{0}$，根据动能定理有

	$3mgh - \mu \cdot 3mg\cos 60^\circ \frac{h}{\sin 60^\circ} = \frac{1}{2} \cdot 3mv_{0}^{2}$

	A 与 B 碰撞后速度分别为 $v_{1}$ 和 $v_{2}$，根据动量守恒、能量守恒有

	$3mv_{0} = 3mv_{1} + mv_{2}$

	$\frac{1}{2} \cdot 3mv_{0}^{2} = \frac{1}{2} \cdot 3mv_{1}^{2} + \frac{1}{2}mv_{2}^{2}$

	B 到达最高点速度为 $v_{3}$，根据牛顿第二定律有

	$mg = m\frac{v_{3}^{2}}{L}$

	根据能量守恒有 $\frac{1}{2}mv_{2}^{2} = \frac{1}{2}mv_{3}^{2} + mg \cdot 2L$

	解得 $v_{2} = \sqrt{3gh}$，$L = 0.6h$。
}

\item
%% number: 19
%% paperName: 2022年海南高考第 19 题 14 分
%% typeId: 6
%% type: 计算题
%% chapter: 磁场
%% point: 安培力
%% method: 无
%% score: 14
%% degree: 350
%% duplicateId: 0
%% body:
如图，水平面上固定有两光滑平行金属导轨，导轨处于磁感应强度为 $B = 0.25 \UT$ 的匀强磁场中，磁场方向竖直向下。金属棒 MN 垂直导轨放置，导轨右侧接自动控制电路，开关 S 接 a 时，电源可使棒中电流大小始终为 $I = 1 \UA$，电流方向可根据需要改变；接 b 时，棒与阻值为 $R = 0.05 \UO$ 的电阻构成回路；电流方向改变及开关切换可瞬间完成。已知棒的质量 $m = 0.1 \Ukg$，电阻 $r = 0.05 \UO$，棒的长度与导轨间距均为 $L = 0.4 \Um$，棒运动过程中始终与导轨接触良好，导轨电路忽略不计且足够长。

\begin{enumerate}
	\item
	当开关 S 接 a，求当棒中电流方向从 M 流向 N 时，导体棒的加速度方向和大小；
	\item
	当开关始终接 a，使棒由静止开始在最短时间内向左运动 $4 \Um$ 后停下，求此过程棒的最大速度；
	\item
	使棒由静止开始在最短时间内向左移动 $7 \Um$ 后停下，求此过程棒中产生的焦耳热。
\end{enumerate}

\onepicture[w=7.50cm, align=right]{figs/19.png}

%% answer:
\jdanswer{
\begin{enumerate}
	\item
	$a = 1 \Umsq$，方向向右

	\item
	$v = 2 \Ums$

	\item
	$Q_{\text{总}} = 0.4 \UJ$
\end{enumerate}
}
%% memo:
\memoanswer{
	（1）当电流从 M 流向 N 时，由左手定则可判断安培力向右，故加速度方向向右。根据牛顿第二定律有 $BIL = ma$

	代入数据可得 $a = 1 \Umsq$

	（2）开关始终接 a 时，电流 N 到 M，经过时间 $t_{1}$ 后电流变为 M 到 N，再经时间 $t_{2}$ 速度减为零，前 $t_{1}$ 时间内有 $x_{1} = \frac{v^{2}}{2a}$

	后 $t_{2}$ 时间内有 $x_{2} = \frac{v^{2}}{2a}$

	根据 $x_{1} + x_{2} = 4$，联立解得 $v = 2 \Ums$

	（3）先接 a 一段时间 $t_{1}$，电流由 N 到 M，再接到 b 端一段时间 $t_{2}$，再接到 a 端一段时间 $t_{3}$，电流由 M 到 N，最后接到 b 静止

	第一段，则有 $v = at_{1}$，$x_{1} = \frac{1}{2}at_{1}^{2}$，$Q_{1} = I^{2}rt_{1}$

	第二段，则由动量定理 $-B\overline{I}Lt_{2} = mv' - mv$

	且 $\overline{I}t_{2} = \frac{BLx_{2}}{R + r}$

	则有 $Q_{2} = \left(\frac{1}{2}mv^{2} - \frac{1}{2}mv'^{2}\right)\frac{r}{R + r}$

	第二段末的加速度与第三段相同，则第三段，

	$\frac{B^{2}L^{2}v'}{R + r} = ma,\quad v' = at_{3},\quad x_{3} = \frac{1}{2}at_{3}^{2},\quad Q_{3} = I^{2}rt_{3}$

	又 $x_{1} + x_{2} + x_{3} = 7$

	解得 $v' = 1 \Ums$，$t_{1} = 3 \Us$，$t_{2} = 1 \Us$

	故 $Q_{\text{总}} = Q_{1} + Q_{2} + Q_{3} = 0.4 \UJ$
}

\end{enumerate}

\end{document}
